\cleardoublepage

\section{结论}

\subsection{总结}

本文围绕基于大语言模型的单元测试生成与修复问题展开研究，针对现有方法在生成阶段和修复阶段普遍存在的单模型依赖、资源利用率不足以及性能提升受限等问题，提出了 HieraTest 方法。该方法通过将任务分解到不同能力层级的模型，并结合静态分析、候选筛选、错误反馈和迭代修复机制，尝试在生成质量和推理成本之间取得平衡。

从论文写作和实验设计的角度看，HieraTest 的核心贡献在于引入了大小模型协同思路，将单元测试生成与测试修复统一纳入一个可迭代的流程中。后续实验将进一步验证该方法在真实开源项目上的有效性，并分析其在不同模型配置下的适用边界。

\subsection{未来工作}

后续工作可从以下几个方向继续推进：一是扩展评估数据集和模型类型，以验证方法的普适性；二是增强错误定位与修复策略，使系统能够更精确地处理复杂编译错误和运行错误；三是进一步优化提示词构造与上下文检索机制，以减少无效调用并提升生成稳定性；四是结合更丰富的测试质量指标，如 mutation score 和 assertion quality，对生成测试进行更全面评估。
